\documentclass[11pt]{amsart}
\usepackage{calc,amssymb,amsmath,ulem,stmaryrd,fullpage}
\usepackage{enumerate}
\usepackage{alltt}
\RequirePackage[dvipsnames,usenames]{color}
\let\mffont=\sf
\normalem
%\input{kmacros3.sty}
\input{xy}
\xyoption{all}
\usepackage{hyperref}
\usepackage{graphicx}
\usepackage[all,cmtip]{xy}
\usepackage{verbatim}
\usepackage[left=0.95in,top=0.95in,right=0.95in,bottom=0.95in]{geometry}
\newcommand{\tauCohomology}{T}
\newcommand{\FCohomology}{S}


%\theoremstyle{theorem}
%\newtheorem*{mainthm*}{Main Theorem}

\newcommand{\note}[1]{\marginpar{\sffamily\tiny #1}}

\def\todo#1{\textcolor{Mahogany}%
{\sffamily\footnotesize\newline{\color{Mahogany}\fbox{\parbox{\textwidth-15pt}{#1}}}\newline}}

\renewcommand{\O}{\mathcal O}
\newcommand{\ie}{{\itshape i.e.} }
\newcommand{\roundup}[1]{\lceil #1 \rceil}
\newcommand{\rounddown}[1]{\lfloor #1 \rfloor}
\renewcommand{\to}[1][]{\xrightarrow{\ #1\ }}
\newcommand{\onto}[1][]{\protect{\xrightarrow{\ #1\ }\hspace{-0.8em}\rightarrow}}
\newcommand{\into}[1][]{\lhook \joinrel \xrightarrow{\ #1\ }}
%\newcommand{DBDual}[1]{{\underline{\omega}_{#1}}}
\newcommand{\DBDual}[1]{{{\uuline \omega}{}^{\mydot}_{#1}}}
\newcommand{\Tt}{{\mathfrak{T}}}
%\newcommand{\bn}{{\mathfrak{n}} }
\newcommand{\sol}[1]{ \vskip 6pt{\color{blue} {\bf Solution:} #1} \vskip 6pt}

\newcommand{\bR}{{\mathbb{R}}}
\newcommand{\va}{{\vec{a}}}
\newcommand{\vb}{{\vec{b}}}
\newcommand{\vzero}{{\vec{0}}}
\newcommand{\vi}{{\vec{i}}}
\newcommand{\vj}{{\vec{j}}}
\newcommand{\vk}{{\vec{k}}}
\newcommand{\vh}{{\vec{h}}}
\newcommand{\vr}{{\vec{r}}}
\newcommand{\vu}{{\vec{u}}}
\newcommand{\vv}{{\vec{v}}}
\newcommand{\vw}{{\vec{w}}}
\newcommand{\vx}{{\vec{x}}}
\newcommand{\vy}{{\vec{y}}}
\newcommand{\vz}{{\vec{z}}}
\newcommand{\vT}{{\vec{T}}}
\newcommand{\vN}{{\vec{N}}}
\newcommand{\vF}{{\vec{F}}}
\newcommand{\vG}{{\vec{G}}}
\newcommand{\bF}{\mathbb{F}}
\newcommand{\bQ}{\mathbb{Q}}
\newcommand{\bZ}{\mathbb{Z}}
\newcommand{\bC}{\mathbb{C}}
\newcommand{\Inn}{\mathrm{Inn}}
\newcommand{\Aut}{\mathrm{Aut}}
\newcommand{\Hom}{\mathrm{Hom}}
\newcommand{\End}{\mathrm{End}}
\newcommand{\coker}{\mathrm{coker}}


\begin{document}
\title{HW \#6 -- Math 6310\\Fall 2019}
\author{Due: Tuesday, November 26th}
\maketitle
\noindent
{\bf 1.}  Suppose $R$ is a commutative ring.  Consider the functor:
\[
    \bullet \mapsto \Hom_R(\Hom_R(\bullet, R), R).
\]
This is called the \emph{reflexification} functor.  

Show that there is a natural transformation of the identity function to the reflexification functor.  (Essentially, I am asking you to show that there is a map from any module $M$ to the reflexification of $M$, and it is compatible with maps from $M \to N$).
\newpage
\noindent
{\bf 2.}  Suppose that $R \to S$ is a map of commutative rings.  Show that $\Hom_R(S, R)$ is an $S$-module.  Then, in the case that $R = \bZ$ and $S = \bZ[i]$, prove that $\Hom_R(S, R)$ is a free $S$-module of rank $1$.  
\newpage
\noindent
{\bf 3.}  Suppose $R$ is a commutative ring.  Given modules $B, \ldots, D, B', \ldots, D'$ and $R$-module maps $b, \ldots, d, a', \ldots, c', \beta, \ldots, \gamma$ as in the commutative diagram below:
\[
\xymatrix{
&  B \ar[r]^b \ar[d]^{\beta} &  C  \ar[r]^c \ar[d]^{\gamma} & D  \ar[d]^{\delta} \ar[r]^d & 0 \\
0 \ar[r]_{a'}        & B' \ar[r]_{b'} & C' \ar[r]_{c'} & D' & \\
}
\]
If the rows are exact, then we have an exact sequence
\[
\ker \beta \to \ker \gamma \to \ker \delta \xrightarrow{\phi} \coker \beta \to \coker \gamma \to \coker \delta
\]
where the maps between the kernels and cokernels are induced by $b,c,b',c'$ and the map $\phi$ is magic (in the proof, the exactness that I really want to read about is the exactness at $\phi$).
\newpage
\noindent
{\bf 4.}  Suppose $R$ is a commutative ring and $I, J \subseteq R$ are ideals.  Prove that
\[
    (R/I) \otimes_R (R/J) \cong R/(I + J)
\]
as $R$-modules.
\newpage
\noindent
{\bf 5.}  Suppose that $M$ is an $R$-module and $W \subseteq R$ is a multiplicative set.  Prove that
\[
    W^{-1} M \cong (W^{-1}R) \otimes_R M
\]
as $R$-modules.  Here $W^{-1} M$ is as defined in 
\newpage
\noindent
{\bf 6.}  A short exact sequence of $R$-modules
\[
    0 \to A \xrightarrow{\alpha} B \xrightarrow{\beta} C \to 0
\]
is called \emph{split} if either of the following equivalent conditions hold.
\begin{itemize}
    \item[(a)]  There exists $f : B \to A$ such that $A \xrightarrow{\alpha} B \xrightarrow{f} A$ is an isomorphism.
    \item[(b)]  There exists $g : C \to B$ such that $C \xrightarrow{g} B \xrightarrow{\beta} C$ is an isomorphism.    
\end{itemize}
Prove they are indeed equivalent.  Then show if the short exact sequence is split, that $B \cong \alpha(A) \oplus C$.
\newpage
\noindent
{\bf 7.}  Suppose that $0 \to A \to B \to C \to 0$ is split.  For any $R$-module $M$, prove that
\[
    \begin{array}{c}
        0 \to \Hom(M, A) \to \Hom(M, B) \to \Hom(M, C) \to 0\\
        0 \to \Hom(C, M) \to \Hom(B, M) \to \Hom(A, M) \to 0\\
        0 \to M \otimes A \to M \otimes B \to M \otimes C \to 0
    \end{array}
\]
are all exact sequences.
\newpage
\noindent
{\bf 8.}  Suppose that $A \to B \to C$ is a sequence of maps between $R$-modules and also that $0 \to \Hom_R(C, N) \xrightarrow{g'} \Hom_R(B, N) \xrightarrow{f'} \Hom_R(A, N)$ is exact for \emph{every} $R$-module $N$.  Prove that $A \to B \to C \to 0$ is exact.
\vskip 3pt
\emph{Hint:} The trick is clever choices of $N$ to get different parts of the exactness.
\vfill
The following is also true:  Suppose that $A \to B \to C$ is a sequence of maps between $R$-modules and also that $0 \to \Hom_R(M, A) \xrightarrow{f''} \Hom_R(M, B) \xrightarrow{g''} \Hom_R(M, C)$ is exact for \emph{every} $R$-module $M$.  Then $0 \to A \to B \to C$ is exact.  I won't ask you to prove it however.
\end{document}

